\documentclass[12pt]{article}
\usepackage[utf8]{inputenc}
\usepackage[T1]{fontenc}
\usepackage{hyperref}
\title{Keoru-Ahia Language Resources for K-12 Curriculum Development}
\author{Compiled by Claude Sonnet 5.5 (claude-sonnet-5-5) for Andr\'as Kornai}
\date{}
\begin{document}
\maketitle

\section{Introduction}
Keoru-Ahia (ISO 639-3 xeu; Glottocode keor1235) is an Eleman language of the Gulf Province, Papua New Guinea, with about 6,000 speakers (2000 census). Its dialects include Keoru, Ahia and Pairi. It is distinct from Kerewo, an unrelated Kiwaian language. Very little material turned up in search.

\section{Primary Educational Resources}
No primers, textbooks or school materials were found.

\section{Linguistic and Reference Materials}
\begin{itemize}
\item Wikipedia article on Keoru: \url{https://en.wikipedia.org/wiki/Keoru_language}
\item Wikipedia article on the Eleman languages: \url{https://en.wikipedia.org/wiki/Eleman_languages}
\item Wiktionary category for Keoru-Ahia (a small collection of entries): \url{https://en.wiktionary.org/wiki/Category:Keoru-Ahia_language}
\item Glottolog language page: \url{https://glottolog.org/resource/languoid/id/keor1235.bigmap.html}
\end{itemize}
No grammar or dictionary was found.

\section{Religious and Community Resources}
Joshua Project reports no complete Bible translation, with a translation begun: \url{https://joshuaproject.net/index.php/languages/xeu} and \url{https://joshuaproject.net/people_groups/12643/PP}. I saw no scripture text URL.

\section{Digital and Audio Resources}
None found.

\section{Curriculum Development Context}
The language is still spoken as a first language across the community, but it appears to be minimally documented and not taught in schools. A curriculum would first need community-led documentation: a wordlist, an orthography and recordings. I found no existing partners or programmes.
\end{document}